\section{Introduction}

A single defining relation can produce very different geometries.
The Baumslag--Solitar groups
\[
 \operatorname{BS}(m,n)
   =\langle a,t\mid ta^mt^{-1}=a^n\rangle,
 \qquad m,n\in\mathbb Z\setminus\{0\},
\]
are the basic obstructions to hyperbolicity among one-relator groups.
Gersten asked whether they are the only subgroup obstructions
\cite[p.~228]{Gersten1992Problems}.
We prove the following statement, resolving Gersten's conjecture positively.
Here $F(X)$ denotes the free group on $X$, and
$\langle X\mid r\rangle$ is its quotient by the normal closure of $r$.

\begin{theorem}\label{thm:main}
Let $X$ be a finite set and let $r\in F(X)$. If the one-relator group
$G=\langle X\mid r\rangle$ contains no subgroup isomorphic to
$\operatorname{BS}(m,n)$ for any nonzero integers $m,n$, then $G$ is
word-hyperbolic.
\end{theorem}

The parameters in this statement may have either sign. In particular,
the hypothesis excludes $\operatorname{BS}(1,1)=\mathbb Z^2$.
We call a group with this exclusion \emph{Baumslag--Solitar-free}.
The customary formulation excluding every $\operatorname{BS}(1,n)$,
for every nonzero integer $n$, is equivalent: every Baumslag--Solitar group contains
one of these groups \cite[Lemma~2.10]{GardamKielakLogan2025Algebraically}.
Theorem~\ref{thm:main} includes trivial relators and proper powers; it
has no torsion-free hypothesis.

We use \emph{free-by-cyclic} in the subgroup-closed sense: a group $H$
fits into an exact sequence
\[
1\longrightarrow K\longrightarrow H\longrightarrow Q\longrightarrow1
\]
with $K$ free of possibly infinite rank and $Q\leq\mathbb Z$, possibly
trivial. Kielak--Linton recall Wise's Conjecture~17.8 that every
hyperbolic one-relator group is virtually free-by-cyclic and prove the
conclusion under the additional hypothesis of virtual compact
specialness \cite[Section~1.1 and Corollary~1.3]{KielakLinton2024Virtually}.
The companion theorem \cite[Theorem~1.1]{OpenAIVirtualSpecialness}
provides that hypothesis. Combining it with Theorem~\ref{thm:main}
gives the Baumslag--Solitar-free form in
Corollary~\ref{cor:free-by-cyclic}, together with a finite-index
embedding into a free-by-$\mathbb Z$ group with finitely generated
free kernel.

\subsection{Earlier work and the remaining step}

The classical theory already separates the torsion case from the
main difficulty. Magnus's Freiheitssatz identifies free subgroups
generated by suitable proper subsets of a one-relator generating set
\cite{Magnus1930}. Newman's spelling theorem for proper-power relators
gives a Dehn algorithm and hence hyperbolicity
\cite[Theorem~3]{Newman1968}; see also the proof in
\cite[Theorem~4]{McCoolSchupp1973}. Thus the new hyperbolicity assertion
concerns non-power relators, although Theorem~\ref{thm:main} includes
all relators. The Freiheitssatz makes an inductive approach possible:
a Magnus splitting has a smaller one-relator vertex group and free
edge groups. The geometric form of the
Magnus--Moldavanskii--Masters hierarchy, developed by Linton, realizes
these splittings through domains in cyclic covers
\cite[Sections~1.1--1.2 and~4]{LintonHierarchies2025}.
We use that construction in Section~\ref{mag:section}, with the domain
and complexity arguments given in the form needed here.

Several recent results identify substantial classes where this
induction preserves hyperbolicity. Louder and Wilton characterized
negative immersions for one-relator presentation complexes by
primitivity rank greater than two
\cite[Theorem~1.3]{LouderWilton2022NegativeImmersions}.
Here the primitivity rank of a relator $r$ is the least rank of a
subgroup of $F(X)$ containing $r$ in which $r$ is not part of a free
basis, with value $\infty$ if no such subgroup exists.
Their subgroup theorem excludes Baumslag--Solitar subgroups in this
class, and their subsequent work establishes uniform negative
immersions and coherence
\cite[Corollary~1.8]{LouderWilton2022NegativeImmersions}
\cite[Theorems~A and~C]{LouderWilton2024UniformNegativeImmersions}.
Linton then proved hyperbolicity and virtual specialness for groups
with negative immersions
\cite[Theorem~7.2]{LintonHierarchies2025}.
His later characterization says that a finitely generated one-relator
group is both hyperbolic and locally quasiconvex precisely when its
relator has primitivity rank different from two
\cite[Theorem~1.4]{Linton2025MappingToriSubgroups}.
Local quasiconvexity means that every finitely generated subgroup is
quasiconvex. This characterization does not settle hyperbolicity
alone for Baumslag--Solitar-free groups whose relator has primitivity
rank two.

Linton's hierarchy criterion makes the remaining obstruction more
specific: a one-relator hierarchy is quasiconvex with hyperbolic
ambient group if and only if it is $\mathbb Z$-stable and the group is
Baumslag--Solitar-free
\cite[Theorem~7.1]{LintonHierarchies2025}.
The stability hypothesis controls the persistence of noncyclic
intersections along the hierarchy; it is an additional condition,
not a consequence supplied by that criterion.
A complementary reduction to primitive-extension groups isolates a
particular class of two-generator one-relator groups in which to
settle the conjecture
\cite[Theorem~5.6 and Corollary~5.7]{LintonHyperbolic2026}.
For a broader historical account, see
\cite{LintonNybergBrodda2025Survey}.

The geometric issue is expansion in a graph of groups. The
Bestvina--Feighn combination theorem makes annular flaring a sufficient
condition for hyperbolicity when the vertex and edge data satisfy the
required geometric hypotheses
\cite[Theorem~1.2]{BestvinaFeighn1996Addendum}.
We use the algebraic annular formulation in
\cite[Theorem~4.5]{Kapovich1999NonQuasiconvexity}.
For mapping tori of free-group automorphisms, Brinkmann proved that
the absence of periodic conjugacy classes forces expansion and
hyperbolicity \cite{Brinkmann2000HyperbolicAutomorphisms}.
Mutanguha extended the hyperbolicity criterion to ascending HNN
extensions of finite-rank free groups: they are hyperbolic exactly
when they contain no Baumslag--Solitar subgroup
\cite[Theorem~6.7]{Mutanguha2021}.
His expanding relative immersions, iterated pullbacks, and annular
arguments supply important precedents for treating persistent
intersections \cite[Theorem~4.5 and Sections~5--6]{Mutanguha2021}.
The combination theorem here permits both free edge subgroups to be
proper in a hyperbolic vertex group, under the intersection and rank
hypotheses below.

\subsection{The combination argument}

Our main technical result is a combination theorem for a restricted
class of HNN extensions of torsion-free hyperbolic groups
(Theorem~\ref{ann:combination}). The Magnus construction reduces the problem to a one-relator vertex
group of smaller complexity and finitely generated free edge groups.
Once the vertex group is hyperbolic, Linton's quasiconvexity theorem gives
quasiconvex edge groups. Two further intersection properties are needed:
conjugates of the edge groups have noncyclic intersection only in one
specified double coset, and their overlap subgroups satisfy a strong
reduced-rank inequality.
Collins's intersection theorem and Linton's strong-inertia result supply
these properties for the Magnus splittings used here
\cite{Collins2008Conjugates,LintonHierarchies2025}.
The combination theorem shows that the absence of Baumslag--Solitar
subgroups forces annular flaring under precisely these hypotheses.

We prove flaring by contradiction, using free-group currents:
invariant measures on pairs of distinct boundary points
\cite{KapovichCurrents2006}.
Long nonflaring annuli give a nonzero stationary current for which
almost every line can be transferred across successive edges of the
Bass--Serre tree.
The stabilizers of increasingly long tree paths are represented by
finite labeled core graphs. The rank inequality bounds the number of
edges after degree-two vertices are suppressed; it does not bound
their label lengths. Along a subsequence, the edges with bounded
labels form one fixed graph, called the small subgraph. Its noncyclic
components determine finitely many maximal subgroup carriers for
persistent path stabilizers. A return along a carrier induces an
atoroidal automorphism, with no nontrivial periodic conjugacy class.
Brinkmann's expansion theorem then forces the set of carrier lines
to have zero current mass.

Zero mass does not exclude a carrier line from the closed support of
the current. This distinction is the main issue in
Section~\ref{exp:section}. Starting with a putative carrier line in
the support, we transport a positive-mass cylinder along expanding
carrier returns.
If this transport continues to follow the prescribed carrier,
recentering and compactness concentrate positive mass on its
zero-mass line set. If it fails, the stabilizers of the competing
actual tree path and the prescribed path have at most cyclic
intersection. A sufficiently long overlap consequently follows a
periodic axis of uniformly bounded primitive period.
Recentering these overlaps puts a periodic line in the support.
We reuse this same failure mechanism to exclude periodic support
lines: pulling the overlap back along the carrier return preserves
enough whole periods to identify the failure axis with the original
periodic trajectory, contradicting the competing path choices.
Counting whole periods is essential because inverse automorphisms
can enlarge primitive period lengths.

It follows that no support line lies in the small subgraph and that
every support line has irrational endpoints. The bounded graph
complexity now leaves only finitely many rays admitting two
different pasts. Transfer permutes their finitely many endpoint
orbits under the free edge group. A return at one of these endpoints
supplies infinitely many disjoint translates of a fixed positive-mass
set inside a compact set,
contradicting local finiteness of the current. This proves the
restricted combination theorem; the geometric Magnus induction
then proves Theorem~\ref{thm:main}.

\subsection{Magnus subgraphs}

We also record the geometric subgroup conclusion in the scope needed
by the induction. A \emph{circuit} in a graph is a map from a subdivided
circle; it is \emph{immersed} when its edge path has no backtracking,
including at the closing point. A nontrivial element is a
\emph{non-power} if it is not $u^k$ with $|k|\geq2$.

\begin{theorem}\label{thm:magnus}
Let $\Gamma$ be a finite connected graph, and let
$\lambda:S^1\to\Gamma$ be a nontrivial immersed circuit representing a
non-power in $\pi_1\Gamma$. Attach a $2$-cell along $\lambda$, obtaining
$Y$. If $\pi_1Y$ is Baumslag--Solitar-free, then, for every connected
subgraph $S\subseteq\Gamma$ with $\lambda(S^1)\nsubseteq S$, the map
\[
 \pi_1S\longrightarrow\pi_1Y
\]
is injective and its image is quasiconvex.
\end{theorem}

A choice of basepoint paths changes these subgroups only by conjugacy.
Such an $S$ is called a \emph{Magnus subgraph}. Its injectivity follows
from the Freiheitssatz. Once Theorem~\ref{thm:main} gives ambient
hyperbolicity, its quasiconvexity follows from Linton's theorem that
Magnus subgroups of hyperbolic one-relator groups are quasiconvex
\cite[Theorem~4.7]{LintonHyperbolic2026}.
We prove the graph-to-basis identification explicitly in
Lemma~\ref{mag:graph-magnus}. The same quasiconvexity theorem supplies the
quasiconvex edge groups at each inductive step, so the induction needs
only to establish hyperbolicity of the smaller vertex group.

\subsection{Structure of the proof}

Sections~\ref{pre:section}--\ref{ray:section} prove the restricted
combination theorem. Section~\ref{ann:section} states its hypotheses and
converts nonflaring annuli into a bounded bilateral sequence of
currents. Section~\ref{stat:section} removes atoms and obtains a
stationary current with well-defined transfers on a conull set.
Sections~\ref{car:section} and~\ref{exp:section} construct the finite
carriers and exclude periodic support lines. Section~\ref{ray:section}
finishes the measure contradiction. Finally, Section~\ref{mag:section}
constructs a decreasing Magnus induction and proves
Theorems~\ref{thm:main} and~\ref{thm:magnus}.
Section~\ref{sec:free-by-cyclic} derives the virtual free-by-cyclic
consequence and records the corresponding ascending-HNN consequence
of the companion theorem.
